\section{Introduction}\label{sec:intro}

A logarithmic-space machine can remember a bounded number of positions
in its input, but cannot keep a polynomial-length computation or random
tape. Can fresh random bits increase what such a machine can decide?
We answer this question by giving a deterministic logarithmic-space
simulation of every polynomial-time randomized logarithmic-space
computation with bounded error.

We use machines with finite control, an endmarked read-only input,
finitely many work tapes, and fresh independent fair bits. Input heads
are confined between the endmarkers. All randomized running-time
bounds hold on every random tape. Space counts traversed writable cells
in bits, including every simultaneously live counter, numerical register
and suspended call; visiting a blank work cell counts toward space.
A transducer has a one-way write-only output tape. Throughout,
$n$ is the input length and logarithms have base two.

The class $\Lclass$ consists of languages decidable by deterministic
machines in $O(\log(n+2))$ work space. Such a halting deterministic
machine has only polynomially many configurations and hence runs in
polynomial time. The class $\RL$ permits
polynomial-time randomized machines in this space, with acceptance
probability zero on no inputs and at least $1/2$ on yes inputs.
For $\BPL$, the corresponding probabilities are at most $1/3$ and at
least $2/3$.

\begin{theorem}[Exact logarithmic-space derandomization]\label{thm:main}
\[
                 \Lclass=\RL=\BPL.
\]
\end{theorem}

The proof approximates arbitrary acceptance probabilities. This gives
the two-sided statement directly, and also applies to computations
whose acceptance probabilities have no decision gap. Its quantitative
form separates the memory needed for the input from the memory needed
for the requested accuracy.

\begin{theorem}[Acceptance-probability approximation]\label{thm:probability}
Fix a randomized machine $\mathcal M$ with polynomial worst-case
running time and $O(\log(n+2))$ work space. Let $p_{\mathcal M}(x)$ be
its acceptance probability on an input $x$ of length $n$. There is a
uniform deterministic algorithm which, on input $(x,1^q)$ with $q\ge1$,
outputs an integer $0\le a\le2^{q+2}$ satisfying
\[
       \left|a2^{-(q+2)}-p_{\mathcal M}(x)\right|\le2^{-q}.
\]
It uses $O(\log(n+2)+q)$ work space and
$(n+2)^{O(1)}2^{O(q)}$ time. The constants depend only on $\mathcal M$.
\end{theorem}

In particular, every fixed inverse-polynomial accuracy is attained in
logarithmic space and polynomial time. The precision parameter is unary;
the running-time bound is exponential in general $q$. We first prove
the fixed inverse-polynomial result and then obtain the single
variable-precision algorithm by padding. Section~\ref{sec:deduction}
also derives total deterministic separators for bounded-error promise
problems. Section~\ref{sec:conclusion} constructs accepting runs when
the acceptance probability has a specified inverse-polynomial lower bound.

The simulation is also effective at the level of machine descriptions.
Given a source program and supplied polynomial-time and logarithmic-space
bounds, one terminating compiler produces a deterministic decider and
explicit numerical bounds for its space and time
(Theorem~\ref{thm:effective-simulation}). The supplied bounds are
hypotheses on the source.
Section~\ref{sec:compiler} constructs this compiler by specializing one
fixed deterministic library.

\subsection{Historical context}

The complexity of probabilistic Turing machines was developed by
Gill~\cite{Gill1977}, including measures of both time and tape usage.
Random walks made the logarithmic-space question concrete:
Aleliunas, Karp, Lipton, Lov\'asz and Rackoff obtained polynomial-time
randomized algorithms for undirected connectivity and explicitly asked
whether polynomial-time randomized logarithmic space could be simulated
in deterministic logarithmic space~\cite{AleliunasKLLR1979}.
Early work also studied probabilistic machines without a polynomial-time
restriction. Borodin, Cook and Pippenger obtained a quadratic-space
deterministic simulation in this broader setting
\cite{BorodinCookPippenger1983}.
The distinction between a space bound and simultaneous time and space
bounds is essential throughout this history.

Pseudorandom generators attack the problem by replacing a long sequence
of fresh random bits with the output generated from a short seed.
Ajtai, Koml\'os and Szemer\'edi gave early logarithmic-space
constructions that hit small-space tests using short random
strings~\cite{AjtaiKomlosSzemeredi1987}. Babai, Nisan and Szegedy
subsequently obtained pseudorandom generators through multiparty
communication bounds~\cite{BabaiNisanSzegedy1992}.
Nisan's generator uses $O(S\log R)$ seed bits for a space-$S$
computation reading $R$ random bits~\cite{Nisan1992}.
For logarithmic space and polynomially many bits this gives
$O(\log^2 n)$ seed length; direct seed enumeration has a
quasipolynomial time bound. Nisan's later simultaneous simulation
achieves polynomial time and $O(\log^2 n)$ space, for two-sided as
well as one-sided bounded error~\cite{Nisan1994RLinSC}.
For $S(n)\ge\log n$, Nisan and Zuckerman showed that a space-$S$
computation using $\operatorname{poly}(S)$ random bits can be simulated
using only $O(S)$ random bits and $O(S)$ space~\cite{NisanZuckerman1996}.
The restriction on the original number of random bits is substantive;
for logarithmic space it permits polylogarithmically many bits.

Saks and Zhou reduced the deterministic space needed for general
bounded-error logarithmic-space computation to
$O(\log^{3/2}n)$~\cite{SaksZhou1999}.
Cai, Chakaravarthy and van Melkebeek established a simultaneous
tradeoff: for each fixed rational $0\le\alpha\le1/2$, the simulation
uses time $n^{O(\log^{1/2-\alpha}n)}$ and space
$O(\log^{3/2+\alpha}n)$
\cite{CaiChakaravarthyVanMelkebeek2006}.
The polynomial-time endpoint has squared-logarithmic space.
Hoza subsequently improved the general space bound to
$O(\log^{3/2}n/\sqrt{\log\log n})$
\cite[Theorem~1.6 and Corollary~A.5]{Hoza2021}; like the bound of
Saks and Zhou, this space containment asserts no polynomial time bound.
The same paper also constructs an improved weighted pseudorandom
generator, in which signed or real weights replace an ordinary
probability distribution; the space bound above does not need this new
generator~\cite[Section~1.5]{Hoza2021}.

Acceptance-probability estimation also connects derandomization to
stochastic matrix products and inverses. Ahmadinejad, Kelner, Murtagh,
Peebles, Sidford and Vadhan developed high-precision small-space
algorithms for random-walk probabilities, with stronger bounds for
Eulerian graphs~\cite{AhmadinejadKelnerMurtaghPeeblesSidfordVadhan2020}.
Cohen, Doron, Sberlo and Ta-Shma later used varying precision across
recursion to improve the space complexity of long stochastic matrix
products~\cite{CohenDoronSberloTaShma2023}. The construction below
works with a finite inverse on a time-layered configuration graph;
its correction and copy hierarchy controls the information retained
while individual entries are evaluated.

Structural restrictions on a computation provide another route.
For undirected connectivity, Nisan, Szemer\'edi and Wigderson obtained
deterministic space $O(\log^{3/2}n)$~\cite{NisanSzemerediWigderson1992},
subsequently improved to $O(\log^{4/3}n)$ by Armoni, Ta-Shma,
Wigderson and Zhou~\cite{ArmoniTaShmaWigdersonZhou2000}.
Reingold's logarithmic-space algorithm settles undirected
connectivity~\cite{Reingold2005,Reingold2008}.
Reingold, Trevisan and Vadhan studied regular directed graphs and gave
logarithmic-space path search for Eulerian directed graphs
\cite{RTV2006}. Their oblivious walk constructions impose additional
labeling conditions. General directed configuration graphs need not
have these properties, so those results did not settle the unrestricted
randomized computation problem.

The relation between one-sided and two-sided derandomization also
requires care.
A read-once branching program is a layered deterministic computation
that reads the input symbols once in a fixed order (bits in the binary
case). Its length is the number of steps,
and its width is the largest number of states in a layer.
A $1/2$-hitting set for a class of programs contains an accepting input
for every member that accepts at least half of all input strings.
Cheng and Hoza proved that one uniformly logarithmic-space-enumerable
family of $1/2$-hitting sets $H_n\subseteq\{0,1\}^n$ for all width-$n$,
length-$n$ read-once branching programs over the binary alphabet would yield
$\mathsf L=\mathsf{BPL}$~\cite[Theorem~1.3]{ChengHoza2022}.
The premise is a hitting-set construction; it is not merely the
language-class equality $\mathsf{RL}=\mathsf L$.

Recent results refine both general and structured models.
Chen, Cohen, Doron, Khaskelberg and Ta-Shma improve error reduction for
weighted generators; in the polynomial-width regime their construction
has seed length $\log|\Sigma|+O(\log^2 T+\log(1/\delta))$, where
$T$ is branching-program length, $\Sigma$ its alphabet and $\delta$
the error~\cite{ChenCohenDoronKhaskelbergTaShma2026}.
Cheng and Wu obtain simultaneous polynomial-time and polylogarithmic-space
simulations for regular branching programs, with extensions to polynomial-time
randomized logarithmic-space machines whose two-way read-only random tape
visits each cell at most constantly often, and to such machines with read-once randomness
and an auxiliary stack whose push, pop and idle schedule is independent
of the random bits. Consecutive stays on a random-tape cell count as one
visit~\cite[Theorems~5--6]{ChengWu2026SC}.
Hardness-versus-randomness approaches additionally give conditional,
certified and per-instance alternatives; Pyne and Tell survey this
line of work~\cite{PyneTell2026}.

Effectivity has its own substantial antecedents.
Pyne, Raz and Zhan construct a universal deterministic estimator whose
$O(S(n))$ space guarantee follows from
$\mathsf{prBPL}\subseteq\mathsf{SPACE}[O(S(n))]$, for
space-constructible $S(n)\ge\log n$
\cite[full version, Theorem~5.1]{PyneRazZhan2023}.
Here $\mathsf{prBPL}$ denotes promise problems solvable in polynomial
time and randomized logarithmic space with bounded error.
The compiler below uses a fixed separator established in this paper
and supplies explicit source-specific numerical resource bounds.
Its construction should be read in this context of universal
simulation, with the source bounds supplied as part of the compiler input.

\subsection{Proof strategy}

The core result is fixed inverse-polynomial approximation
(Theorem~\ref{thm:fixed-precision}). Fix the source machine and the
accuracy exponent, and put $B=\Theta(\log(n+2))$. We construct a finite
list of estimates indexed by $O(B)$-bit descriptions. More than three
quarters are accurate, and every estimate can be evaluated by a halting
deterministic computation in $O(B)$ space. Their median gives the desired
approximation. The descriptions encode the possible random
\emph{environments}, of which there are at most $2^{O(B)}$. An environment
consists of a fixed number of \emph{sampling matrices}
(Section~\ref{sec:env}): random matrices over a prime field of size
$2^{O(B)}$, conditioned on their positive-degree coefficient columns
being linearly independent. Applied to the vector of powers of a vertex
identifier, each matrix decides whether that vertex is \emph{retained}
at a prescribed sampling rate. A vertex is retained with probability
comparable to the rate, and, conditioned on the retention of one
identifier, a distinct identifier is retained with probability at most
a constant times its rate (Lemma~\ref{lem:env-ranks}). The environment
also contains $O(B)$ independent fair bits, used as fingerprint seeds
(Section~\ref{sec:fp-blocks}). The median list uses only the encodings
that retain the designated start of the final reward query; it is under
this conditional law that the stated majority is accurate
(Section~\ref{sec:conclusion}). Each estimator task carries an integer
\emph{statistical budget}~$M$. Its error bound in the metrics of
Section~\ref{sec:atom} is $C_b2^{-p}$ for a fixed constant $C_b$, with
precision $p$ growing linearly in $M$
(Equation~\eqref{eq:est-precision}). Every other
estimator task it queries has a smaller budget. The top-level task has
budget $M_{\rm top}=O(B)$. The proof must control both estimation error in a shared
random environment and the storage retained across recursive calls;
enumeration requires termination even on inaccurate choices.

Two resource costs in the preceding literature help locate this approach.
Direct enumeration of Nisan's generator uses an $O(\log^2 n)$-bit seed
\cite{Nisan1992}. Saks and Zhou's matrix-powering algorithm, specialized
to logarithmic space, has $O(\sqrt{\log n})$ recursive levels using
$O(\log n)$ space each \cite[Section~1, pp.~377--379]{SaksZhou1999}.
Here we enumerate an $O(\log n)$-bit environment for an input-specific
probability estimator. Reusing that environment makes the recursive
errors dependent, while keeping a full endpoint or numerical word at
each level would exceed the desired space bound. The construction below
addresses these costs together: conditional averaging and the block
comparison control the shared errors; compression, outgoing and incoming port labels and
allowance differences control the retained storage. This is a direct
estimation argument.

\paragraph{Transferring transitions into rewards.}
Section~\ref{sec:model} represents the computation by a substochastic
matrix $S$ on time-labelled configurations $V_0$, with final accepting
indicator $e$. Every transition advances time, so
$p_0=(I-S)^{-1}e$ is a finite sum and records acceptance probabilities
at all configurations. Section~\ref{sec:hierarchy} successively replaces
transitions by accumulated rewards. For a current transition matrix $C$,
the identity
\[
 D=C-E+EC,\qquad I-D=(I+E)(I-C)
\]
allows a nonnegative correction $E$ to update a reward $W$ to $(I+E)W$
while preserving the represented probability vector. The aim is to make
the remaining transition contribution small.

Correction and copying serve complementary purposes. A finite detector
calculation suppresses correction in columns with many substantial
entries. Where fully enabled, correction makes entries small. The
remaining large entries are divided among copies of their destinations. The detector
bounds how many vertices need active extra copies, so each stage reduces
the entry bound by $2^H$ while at most doubling the number of vertices
incident to positive entries, for fixed $H\ge4$
(Proposition~\ref{prop:hier-copy-step}). Correction itself needs only
constantly many iterations: at its comparison fixed point, a positive
lower bound on the residual at corrected entries bounds their number
in each row. Strictly increasing paths through that support give the
factorial truncation bound of Lemma~\ref{lem:hier-fixed-point}.
The remaining transition mass in each row is at most the entry bound
times the number of active vertices. Transporting rewards through the
copy projections therefore yields accumulated rewards $W_l$ with
\[
 0\le p_0(x_0)-W_l(z_l)\le |V_0|\,2^{-(H-1)l},
\]
where $x_0$ is the start and $z_l$ its copy obtained by choosing the
designated copy labelled zero at each stage. Thus $L=O(B)$ stages suffice;
the full stage domains still have polynomial size. This is the target
proved in Proposition~\ref{prop:hier-reward-target}.

\paragraph{Estimating in one shared environment.}
The hierarchy's exact matrices and conditional means define targets for
recursive queries. Section~\ref{sec:atom} represents their contributions
by random tables on the retained vertex sets of Section~\ref{sec:env}.
Cutoffs suppress crowded sampled rows and, for corrections, crowded
columns. On the support of a true contribution, the conditional mean
of its cutoffs is bounded away from zero, so division by that mean
recovers the contribution in expectation. Clipped denominators keep
values bounded, and the bounds on row support and incoming degree hold
even when estimated inputs are inaccurate. Substitution of earlier
estimates introduces bias, which
Section~\ref{sec:est} bounds.

The incoming bound connects the error and space arguments. After a
correction move, the algorithm can discard the old source and query a
right factor at the intermediate vertex, retaining only a bounded
inverse port for the return. In the error analysis, the same incoming
bound lets us sum out the old source, at a cost given by the ratio of
sampling rates.
Lemma~\ref{lem:est-mean-comparison} gives this cost, and the precision
assignment in Section~\ref{sec:est} pays it.

To approximate the conditional means, Section~\ref{sec:averaging}
constructs walks preserving the specified source and endpoint retention
tests. It derives a uniform spectral gap for the required finite affine
actions from Shalom's property~(T) results
\cite[Theorems~2.1 and~2.6]{Shalom1999}. The resulting conditional
averaging contracts eighth-moment noise exponentially in walk length.
Its steps can be described independently of the conditioning vertices,
with validity checked when those vertices are visited
(Lemma~\ref{lem:env-variants}); this permits their use with one current
vertex address.

\paragraph{Differences and row compression.}
At parent budget $M$, the additive schedule in Section~\ref{sec:est}
averages differences between child estimates at budgets $m$ and $m-1$.
Both approximate the same sample term formed from the true tables, so
their difference has small absolute norm. Their conditional means telescope, and a walk
whose length is proportional to the slack $M-m$ reduces the remaining
noise. Contributions at the same endpoint are combined before taking
the difference. The same slack bounds the immediate path description,
linking the averaging available to the storage it requires.

Section~\ref{sec:buf-compression} groups comparable slacks and represents
each row by a finite list of exceptions to one common default. Small
entries are replaced by that default. At the first overflow, where
capacity forces the threshold above its prescribed floor, either many
retained entries are truly large, an event rare even conditional on a
specified endpoint, or the prefix errors account for the excess. The analysis
uses this rarity for the deterministic conditional-mean error and the
absolute moment bound for the noise, which may depend on the overflow
event. Stability of the later updates then controls the bounded row
functions needed by the hierarchy (Proposition~\ref{prop:buf-error}).

\paragraph{Fingerprints for adaptive endpoints.}
Even these lists cannot be evaluated by saving a full endpoint at every
recursive level. Section~\ref{sec:fp-blocks} uses short fingerprints
at small slacks, combining polynomial evaluation and affine
hashing~\cite{Schwartz1980,CarterWegman1979}. With top budget $O(B)$,
sharing one $O(\log B)$-bit seed over each block of $\Theta(\log B)$
budgets uses $O(B)$ bits. A fixed-pair collision bound is insufficient
because earlier comparisons affect later candidates.

The analysis therefore replaces fingerprint comparisons by exact
equality throughout the current block, retaining the actual procedures
below it. This comparison calculation supplies a small set of candidate
endpoints independent of the current fingerprint. An actual path
escaping that set contains an edge where a smooth support function is
one and its counterpart in the comparison calculation is zero. The
lower-call error bound pays for such escapes; inside the set the
ordinary collision bound applies. Only new increments are combined by
fingerprint: old values stay at their own endpoints. Thus the collision
cost is proportional to the small consecutive-budget differences, rather than
to accumulated values. Together with the compression estimate, this
coupled induction proves the moment bound for the actual estimator in
Theorem~\ref{thm:fp-estimators}.

\paragraph{Total evaluation in logarithmic space.}
Sections~\ref{sec:arith} and~\ref{sec:ctrl} implement these estimates
using one vertex cursor and numerical queries returning one bounded
digit at a time. Lists and ports are fixed independently of the
requested digit, so greater numerical precision approximates the same
estimator value. Incoming accesses use the forest traversal of Cook
and McKenzie~\cite[Proposition~1]{CookMcKenzie1987}, also used in the
reversible simulation of Lange, McKenzie and Tapp~\cite{LMT2000}.
Here the traversed tree is the component of a terminal program's
success state; its canonical initial states enumerate the forward
ports ending there. Returning to that success state recovers the
endpoint without a saved address (Section~\ref{sec:ctrl-incoming}).

Even a short fingerprint index would cost too much if saved at every
call with small budget slack. Lemma~\ref{lem:fp-catalytic} compares two fingerprints by toggling
and reading a shared $B$-bit vector on repeated endpoint visits. The
initial contents cancel and are restored. This adds only constant local
control beyond the endpoint descriptions, and nested comparisons reuse
the same vector. This vector is counted in the work space.

The remaining records are charged to differences between caller and
child resources. If their recursion allowances are $w,w'$ and their
digit places are $\nu,j$, the suspended caller uses at most
\[
                    A_{\rm space}(w-w')+D_{\rm space}(\nu-j)
\]
bits for fixed $A_{\rm space},D_{\rm space}$, including local copies used by a traversal.
The top allowance and digit place are $O(B)$. Allowance strictly
decreases; digit place can increase only in proportion to that decrease.
These differences telescope along the active call chain, while the
environment, cursor, shared bit vector and scratch occur once.
Theorem~\ref{thm:ctrl-realization} proves the resulting
$O(B)$ bound and termination for every environment.

\paragraph{The median and the other results.}
The moment bound is integrated over vertices. Conditioning on retention
of the designated start loses at most a constant times the stage-size
bound $N=2^{C_NB}$, for fixed $C_N$. Increasing the statistical budget by
a fixed multiple of $B$ pays this loss. Section~\ref{sec:conclusion}
combines this with the hierarchy and digit errors to obtain the promised
majority of accurate dyadic estimates. It enumerates the original
environment encodings with their multiplicities, preserving exactly
the conditional law. Repeated enumeration finds the median with only
logarithmic counting and numerical registers. There are polynomially
many encodings, and configuration counting bounds every total
logarithmic-space digit query by polynomial time. This proves
Theorem~\ref{thm:fixed-precision}; constant accuracy and threshold $1/2$
give Theorem~\ref{thm:main} and total promise separators in
Section~\ref{sec:deduction}.

For the unary precision request in Theorem~\ref{thm:probability},
increasing digit accuracy alone leaves the statistical error unchanged.
Section~\ref{sec:deduction} instead fixes an inverse-linear approximator
for a machine that ignores padding. A virtual padded input long enough
for that error to meet the request preserves the original acceptance
probability. Generating each requested input bit from its address gives
the stated $O(\log(n+2)+q)$ space and
$(n+2)^{O(1)}2^{O(q)}$ time. For accepting computations,
Corollary~\ref{cor:certified-search} repeatedly chooses the successor
with largest estimated continuation probability. The total estimation
loss is smaller than the promised lower bound on acceptance, forcing
an accepting terminal state. A verification pass followed by replay
ensures that only accepting computations are output, including outside
that promise.

Finally, the compiler in Section~\ref{sec:compiler} fixes one universal
randomized interpreter that enforces resource caps on every input.
Probability approximation supplies a total deterministic separator for
this interpreter even though arbitrary encoded inputs need not satisfy
a bounded-error promise. The supplied source bounds determine padding
that makes the caps harmless and preserves acceptance probability
exactly. Specializing the fixed separator's virtual input operations
produces the required machine. Its library and input routines give
explicit space bounds; counting complete configurations of that same
machine gives explicit time bounds (Section~\ref{subsec:compiler-bounds}).
